% ===================================================================
%  Dodatek 0: Aksjomatyka różnicy — Z_2 jako akt ustanowienia substratu
%  TGP v1 — Warstwa pre-substratowa (przed Γ, przed wymiarem, przed wszystkim)
% ===================================================================
\section{Aksjomatyka r\'o\.znicy: $\Zp$ jako akt ustanowienia substratu}%
\label{app:aksjomatyka-roznicy}

Niniejszy dodatek formalizuje \textbf{warstw\k{e} pre-substratow\k{a}}
TGP --- logicznie wcze\'sniejsz\k{a} ni\.z substrat $\Gamma$, wymiar
$d=3$, a~tym bardziej ni\.z pole $\Phi$. Proponuje si\k{e} promocj\k{e}
dotychczasowej \emph{propozycji} o~symetrii $\Zp$
(\ref{prop:Z2} w~\S\ref{ssec:chiralna}) do rangi \textbf{Aksjomatu~0}:
symetria $\Zp$ nie jest cech\k{a} hamiltonianu substratu, lecz
\textbf{aktem ustanowienia samego substratu}.

Intuicja przewodnia: stan nicości $\Nzero$ nie jest
niestabilny ,,fizycznie'' (bo fizyka wymaga ju\.z substratu), lecz
\textbf{niestabilny logicznie} --- nie istnieje niezdegenerowany
opis $\Nzero$ bez wprowadzenia pierwszej r\'o\.znicy, a~najmniejsza
nietrywialna r\'o\.znica to w\l{}a\'snie $\Zp$. Substrat $\Gamma$
wy\l{}ania si\k{e} w konsekwencji jako \emph{minimalna realizacja
konceptu r\'o\.znicy}.

Taki uk\l{}ad zamyka \l{}a\'ncuch przyczynowy:
\[
  \Nzero \;\longrightarrow\; \Zp
  \;\longrightarrow\; \hat{s}_i \in \{+,-\}
  \;\longrightarrow\; \Gamma
  \;\longrightarrow\; \Phi
  \;\longrightarrow\; g_{\mu\nu}
  \;\longrightarrow\; \text{SM + kosmologia},
\]
przy czym ka\.zda strza\l{}ka jest albo twierdzeniem ([AN]), albo
konieczno\'sci\k{a} logiczn\k{a} ([LOG]), z~dok\l{}adnie jedn\k{a}
hipotez\k{a} [HYP] (CV=1, por.\ dod.~\ref{app:R2}).

% -------------------------------------------------------------------
\subsection{Motywacja: dlaczego ,,pod substratem'' nie jest dobrze
            postawionym pytaniem}%
\label{app:0-motywacja}
% -------------------------------------------------------------------

Fizyka operuje na \emph{r\'o\.znicach}. Obserwable to zawsze relacje
mi\k{e}dzy stanami, nie ,,stany bezwzgl\k{e}dne''. St\k{a}d pytanie
,,co jest pod $\Gamma$?'' wymaga punktu odniesienia
\emph{spoza} $\Gamma$, ale je\'sli $\Gamma$ jest fundamentalny, to
takiego punktu z~definicji nie ma --- pytanie unieważnia si\k{e}
operacyjnie.

Pozostaje jednak \textbf{sensowne pytanie przeformu\l{}owane}: czy
$\Gamma$ jest \emph{konieczny} w~sensie strukturalnym, tzn.~czy
istnieje minimalna forma matematyczna, z~kt\'orej $\Gamma$ wy\l{}ania
si\k{e} bez dodatkowego wyboru? Celem niniejszego dodatku jest
odpowied\'z twierdz\k{a}ca: t\k{a} form\k{a} jest \textbf{koncept
r\'o\.znicy}, a~jego minimaln\k{a} matematyczn\k{a} realizacj\k{a}
--- grupa $\Zp$.

\begin{remark}[Horyzont epistemiczny]
\label{rem:0-horyzont}
$\Gamma$ nie jest ,,\l{}ukiem'' TGP, lecz jej \textbf{horyzontem
epistemicznym}. Aksjomat~0 nie usuwa tego horyzontu --- przesuwa go
z~,,substrat jest postulowany'' na ,,r\'o\.znica jest postulowana''.
R\'o\.znica zaś jest warunkiem mo\.zliwo\'sci jakiegokolwiek
dyskursu (tak\.ze dyskursu o~nicości), wi\k{e}c jest to
\emph{najtańszy mo\.zliwy} punkt oparcia.
\end{remark}

% -------------------------------------------------------------------
\subsection{Aksjomat~0 i jego trzy r\'ownowa\.zne sformu\l{}owania}%
\label{app:0-aksjomat}
% -------------------------------------------------------------------

\begin{axiom}[R\'o\.znica jako warunek istnienia]\label{ax:0-roznica}
Istnienie czegokolwiek wy\l{}\k{a}cza nicość absolutn\k{a} $\Nzero$.
Najmniejsza nietrywialna struktura zgodna z~tym wy\l{}\k{a}czeniem
to dyskretna, inwolutywna grupa rz\k{e}du~2:
\[
  \Zp \;=\; \{\mathbf{1},\,\sigma\},
  \qquad \sigma^2 = \mathbf{1},
  \qquad \sigma \neq \mathbf{1}.
\]
Generator $\sigma$ nazywamy \textbf{aktem r\'o\.znicy}; jego
dzia\l{}anie na dowoln\k{a} struktur\k{e} pochodn\k{a} indukuje
chiraln\k{a} inwersj\k{e} $\hat{s} \mapsto -\hat{s}$.
\end{axiom}

Aksjomat~\ref{ax:0-roznica} ma trzy r\'ownowa\.zne sformu\l{}owania
--- kategoryjne, algebraiczne i dynamiczne --- kt\'ore poni\.zej
przedstawiamy jako twierdzenia o~r\'ownowa\.zności (dowody w~app.~\ref{app:0-dowody}).

% -------------------------------------------------------------------
\subsubsection{Sformu\l{}owanie kategoryjne}%
\label{app:0-kategoryjne}
% -------------------------------------------------------------------

\begin{proposition}[Endofunktor negacji na kategorii pustej]%
\label{prop:0-endofunktor}
Niech $\mathbf{Cat}_{\min}$ oznacza klas\k{e} kategorii posiadaj\k{a}cych
endofunktor $\neg$ spe\l{}niaj\k{a}cy $\neg\circ\neg = \mathrm{id}$ i
$\neg \neq \mathrm{id}$. W\'owczas:
\begin{enumerate}[label=(\roman*), nosep]
  \item Kategoria pusta $\mathbf{0}$ nie nale\.zy do $\mathbf{Cat}_{\min}$
    (endofunktor $\neg$ wymaga dziedziny niepustej).
  \item Minimalny obiekt $\mathbf{Cat}_{\min}$ to kategoria dyskretna
    z~dok\l{}adnie dwoma obiektami $\{+, -\}$ i~dwiema identyczno\'sciami,
    z~$\neg(+) = -$, $\neg(-) = +$.
  \item Grupa automorfizm\'ow tej kategorii jest izomorficzna z~$\Zp$.
\end{enumerate}
\end{proposition}

Interpretacja: sam \emph{koncept negacji} --- rozumiany jako
endofunktor inwolutywny, nietrywialny --- wymusza istnienie co
najmniej jednego obiektu (bo inaczej funktor nie ma na czym
dzia\l{}a\'c) i~grupy jego automorfizm\'ow, kt\'ora jest dok\l{}adnie
$\Zp$. Nicość $\mathbf{0}$ jest w~tym sensie \textbf{niestabilna
kategoryjnie}: nie przyjmuje nietrywialnego endofunktora
inwolutywnego.

% -------------------------------------------------------------------
\subsubsection{Sformu\l{}owanie algebraiczne (Clifford)}%
\label{app:0-clifford}
% -------------------------------------------------------------------

\begin{proposition}[Minimalna wie\.za Clifforda]\label{prop:0-clifford}
Rozwa\.zmy rodzin\k{e} algebr Clifforda $\mathrm{Cl}(p,q)$ nad
$\mathbb{R}$. W\'owczas:
\begin{enumerate}[label=(\roman*), nosep]
  \item $\mathrm{Cl}(0,0) \cong \mathbb{R}$ --- algebra trywialna,
    brak struktury rozr\'o\.zniaj\k{a}cej.
  \item $\mathrm{Cl}(1,0) \cong \mathbb{R} \oplus \mathbb{R}$
    --- pierwsza algebra z~nietrywialn\k{a} $\Zp$-gradacj\k{a}.
  \item Gradacja $\Zp$ w~$\mathrm{Cl}(1,0)$ jest \textbf{jedyn\k{a}
    mo\.zliw\k{a}} gradacj\k{a} nie wymagaj\k{a}c\k{a} uprzedniego
    za\l{}o\.zenia o~ciele liczb zespolonych ani o~wymiarze.
\end{enumerate}
\end{proposition}

Interpretacja: przej\'scie od ,,braku struktury'' do ,,minimalnej
struktury z~rozr\'o\.znieniem'' polega na dodaniu jednego generatora
$e$ z~$e^2 = +1$, co automatycznie wprowadza $\Zp$-gradacj\k{e}.
Dalsze wie\.ze Clifforda (spin, orientacja, geometria spinorowa)
wy\l{}aniaj\k{a} si\k{e} jako iteracje tego samego aktu.

% -------------------------------------------------------------------
\subsubsection{Sformu\l{}owanie dynamiczne (niestabilność degeneracji)}%
\label{app:0-dynamiczne}
% -------------------------------------------------------------------

\begin{proposition}[Niestabilność punktu sta\l{}ego identyczno\'sci]%
\label{prop:0-niestabilnosc}
Niech $\mathcal{T}$ b\k{e}dzie przestrzeni\k{a} stan\'ow, a~$\mathrm{id}:
\mathcal{T} \to \mathcal{T}$ operatorem identyczno\'sciowym. W\'owczas:
\begin{enumerate}[label=(\roman*), nosep]
  \item Ka\.zdy stan $x \in \mathcal{T}$ jest punktem sta\l{}ym
    $\mathrm{id}$ --- ca\l{}a $\mathcal{T}$ jest zdegenerowana.
  \item Zdegenerowany zbi\'or punkt\'ow sta\l{}ych nie wyznacza
    \.zadnego stanu wyr\'o\.znionego; operacyjnie r\'ownowa\.zny
    nicości $\Nzero$.
  \item Pierwsza nieidentyczno\'sciowa inwolucja $\sigma: \mathcal{T}
    \to \mathcal{T}$ z~$\sigma^2 = \mathrm{id}$, $\sigma \neq \mathrm{id}$
    roz\l{}amuje t\k{e} degeneracj\k{e} na dwie orbity $\{x, \sigma x\}$,
    generuj\k{a}c $\Zp$-strukturę.
\end{enumerate}
\end{proposition}

Interpretacja: ,,nicość'' rozumiana jako stan o~maksymalnej
degeneracji jest dynamicznie niestabilna --- najmniejsze zaburzenie,
kt\'ore roz\l{}amuje degeneracj\k{e}, musi by\'c inwolucj\k{a}
(bo inaczej nie powr\'oci do stanu wyj\'sciowego), a~najmniejsza
nietrywialna inwolucja generuje $\Zp$.

\begin{theorem}[R\'ownowa\.zność trzech sformu\l{}owa\'n]%
\label{thm:0-rownowaznosc}
Sformu\l{}owania~\ref{prop:0-endofunktor},
\ref{prop:0-clifford} i~\ref{prop:0-niestabilnosc} s\k{a}
logicznie r\'ownowa\.zne w~nast\k{e}puj\k{a}cym sensie: ka\.zde z~nich
wymusza istnienie minimalnej struktury inwolutywnej
$(\mathcal{T}_{\min}, \sigma)$ z~$|\mathcal{T}_{\min}| = 2$ i~grup\k{a}
automorfizm\'ow $\Zp$, a~ka\.zda dw\'ojka takich struktur jest
izomorficzna.
\end{theorem}

% -------------------------------------------------------------------
\subsection{\L{}a\'ncuch przyczynowy: od $\Nzero$ do $\Gamma$}%
\label{app:0-lancuch}
% -------------------------------------------------------------------

Aksjomat~\ref{ax:0-roznica} zamyka \l{}a\'ncuch przyczynowy TGP
od strony ,,wej\'sciowej''. Poni\.zej formalna rekonstrukcja.

\begin{chain}[Krok 0.1: od $\Nzero$ do dw\'ojki $\{+,-\}$]%
\label{chain:0-krok1}
Aksjomat~\ref{ax:0-roznica} $\Rightarrow$
istnieje co najmniej jeden obiekt i~jego $\Zp$-obraz.
Zbi\'or minimalny: $\{+, -\}$. Status: \textbf{[LOG]}.
\end{chain}

\begin{chain}[Krok 0.2: od $\{+,-\}$ do sieci $V$]%
\label{chain:0-krok2}
Jeden obiekt z~$\Zp$-inwolucj\k{a} nie mo\.ze generowa\'c dynamiki
(brak s\k{a}siad\'ow, brak sprz\k{e}\.ze\'n). Minimalna struktura
relacyjna to \emph{zbi\'or} $V$ z~dyskretn\k{a} topologi\k{a}
s\k{a}siedztwa. \textbf{Liczba element\'ow $V$ nie jest jednak
okre\'slona przez aksjomat~\ref{ax:0-roznica}}; jest to dodatkowy
stopień swobody, kt\'ory domyka aksjomat relacyjno\'sci
(A1, \S\ref{ssec:aksjomaty}). Status: \textbf{[LOG + aks.~A1]}.
\end{chain}

\begin{chain}[Krok 0.3: od sieci $V$ do hamiltonianu $H_\Gamma$]%
\label{chain:0-krok3}
Wymaganie symetrii $\Zp$ ka\.zdemu wierzcho\l{}kowi $i \in V$ przypisuje
zmienn\k{a} $\hat{s}_i \in \mathbb{R}$ z~transformacj\k{a}
$\hat{s}_i \mapsto -\hat{s}_i$.
Najbardziej og\'olny $\Zp$-niezmienniczy hamiltonian lokalny z~czynnikiem
sprz\k{e}\.zenia do najbli\.zszych s\k{a}siad\'ow to:
\begin{equation}\label{eq:0-H-general}
  H_\Gamma = \sum_i \Bigl[
    \tfrac{1}{2}\hat{\pi}_i^2
    + \tfrac{m_0^2}{2}\hat{s}_i^2
    + \tfrac{\lambda_0}{4}\hat{s}_i^4
    + \mathcal{O}(\hat{s}_i^6)
  \Bigr]
  - J\!\sum_{\langle ij\rangle} \hat{s}_i\hat{s}_j.
\end{equation}
Wy\.zsze pot\k{e}gi $\hat{s}_i^{2n}$ s\k{a} \emph{dozwolone}, lecz
\textbf{nieistotne w~IR} (uniwersalno\'s\'c Wilsona--Fishera,
dod.~\ref{app:N}). Status: \textbf{[LOG + uniwersalno\'s\'c IR]}.
\end{chain}

\begin{chain}[Krok 0.4: od $H_\Gamma$ do $\Phi$]%
\label{chain:0-krok4}
Coarse-graining $\Phi(x) = \langle \hat{s}^2 \rangle_{\mathrm{blok}}$
(dod.~\ref{app:substrat}, r\'own.~\ref{eq:B-block}) jest
\textbf{$\Zp$-niezmienniczy} ($\hat{s}^2 \to \hat{s}^2$), wi\k{e}c
$\Phi \geq 0$ automatycznie (aksjomat~\ref{ax:przestrzen}).
Status: \textbf{[AN]}.
\end{chain}

\begin{chain}[Krok 0.5: od $\Phi$ do geometrii]%
\label{chain:0-krok5}
Metryka efektywna $g_{\mu\nu}^{\mathrm{eff}}$ i sta\l{}e $c(\Phi)$,
$\hbar(\Phi)$, $G(\Phi)$ wy\l{}aniaj\k{a} si\k{e} jako funkcjona\l{}y
gradient\'ow $\Phi$ (sek.~\ref{sec:formalizm}). Status: \textbf{[AN]}.
\end{chain}

\begin{corollary}[Ontologiczna jednoznaczno\'s\'c]\label{cor:0-jednoznacznosc}
Przy za\l{}o\.zeniu aksjomatu~\ref{ax:0-roznica}, aksjomatu
relacyjności A1 oraz wymiaru $d=3$ (dod.~\ref{app:R2}, krok
\emph{wymiar}), struktura TGP jest \textbf{jednoznaczna w~IR}
z~dok\l{}adnością do klasy uniwersalności 3D Isinga.
\end{corollary}

% -------------------------------------------------------------------
\subsection{Relacja do dotychczasowych aksjomat\'ow TGP}%
\label{app:0-relacja}
% -------------------------------------------------------------------

Dotychczasowa struktura aksjomatyczna TGP (sek.~\ref{sec:ontologia})
sk\l{}ada\l{}a si\k{e} z:
\begin{itemize}[nosep]
  \item \textbf{A1} --- relacyjno\'s\'c substratu,
  \item \textbf{A2} --- d\'zwigar ($\Phi \geq 0$),
  \item \textbf{A3} --- dynamika (Hamilton/Lagrange),
  \item \textbf{A4} --- niestabilność $\Nzero$,
  \item \textbf{P1--P3} --- postulaty pomocnicze,
  \item \textbf{prop:Z2} --- symetria $\Zp$ jako w\l{}asność
    $H_\Gamma$.
\end{itemize}

Niniejszy dodatek proponuje \textbf{reorganizacj\k{e}}:

\begin{table}[h]
\centering
\begin{tabular}{lll}
\toprule
Status dotychczasowy & Status proponowany & Uzasadnienie \\
\midrule
prop:Z2 (w\l{}asno\'s\'c $H_\Gamma$) &
  \textbf{Aksjomat 0} (konstytutywny) &
  $\Zp$ \emph{ustanawia} $\Gamma$, nie jest \emph{cech\k{a}}
  gotowego $\Gamma$ \\
A4 (niestabilność $\Nzero$) &
  \emph{Twierdzenie} z~aks.~0 &
  Wynika z~prop.~\ref{prop:0-niestabilnosc} \\
A2 ($\Phi \geq 0$) &
  \emph{Twierdzenie} z~aks.~0 &
  Wynika z~$\Zp$-parzysto\'sci $\hat{s}^2$ \\
A1 (relacyjno\'s\'c) &
  \textbf{Aksjomat~1} (bez zmian) &
  Niezale\.zny stopień swobody: \emph{rozmiar} $|V|$ \\
A3 (dynamika) &
  \textbf{Aksjomat~2} (bez zmian) &
  Niezale\.zny: wyb\'or re\.zimu czasowego $\tau$ \\
\bottomrule
\end{tabular}
\caption{Propozycja reorganizacji bazy aksjomatycznej TGP po
promocji $\Zp$ do Aksjomatu~0.}
\label{tab:0-reorganizacja}
\end{table}

\textbf{Redukcja:} baza aksjomatyczna maleje z~4+3 = 7 element\'ow
do 3 aksjomat\'ow (0, 1, 2) + aksjomat wymiaru ($d=3$, dod.~\ref{app:R2})
+ jedna hipoteza CV=1 (dod.~\ref{app:R2}). To istotne zag\k{e}szczenie
strukturalne bez utraty mocy predykcyjnej.

% -------------------------------------------------------------------
\subsection{Co aksjomat~0 domyka, a~co pozostawia otwarte}%
\label{app:0-otwarte}
% -------------------------------------------------------------------

\paragraph{Domyka:}
\begin{itemize}[nosep]
  \item pytanie ,,dlaczego $\Zp$ a~nie inna grupa'' ---
    $\Zp$ jest minimalna logicznie (tw.~\ref{thm:0-rownowaznosc});
  \item pytanie ,,sk\k{a}d $\Phi \geq 0$'' ---
    z~$\Zp$-parzysto\'sci (krok~\ref{chain:0-krok4});
  \item pytanie ,,dlaczego niestabilność $\Nzero$'' ---
    z~niestabilności degeneracji (prop.~\ref{prop:0-niestabilnosc});
  \item pytanie ,,dlaczego kwartyczny $\phi^4$'' ---
    nie dos\l{}ownie, ale w~IR przez uniwersalno\'s\'c
    (krok~\ref{chain:0-krok3}).
\end{itemize}

\paragraph{Nie domyka (uczciwie):}
\begin{itemize}[nosep]
  \item \textbf{Rozmiar $|V|$ i topologia grafu $\Gamma$}: aksjomat~0
    wymusza istnienie $V$, ale nie jego kardynalno\'s\'c ani struktur\k{e}
    s\k{a}siedztwa. To jest w\l{}a\'snie rol\k{a} aksjomatu~A1.
  \item \textbf{Wymiar $d=3$}: nie wynika z~$\Zp$; pozostaje zamkni\k{e}ty
    niezale\.znie przez zbie\.zno\'s\'c ogona solitonu (dod.~\ref{app:R2}).
  \item \textbf{Hipoteza CV=1}: $Q_K = 3/2$ wymaga trzeciego ogniwa
    w~\l{}a\'ncuchu entropijnym --- aksjomat~0 nie wp\l{}ywa na t\k{e}
    spraw\k{e}.
  \item \textbf{Istnienie ,,konceptu r\'o\.znicy''}: pytanie
    \emph{dlaczego w~og\'ole negacja}. Moim zdaniem to pytanie le\.zy
    poza fizyk\k{a} i~poza matematyk\k{a} --- jest warunkiem
    mo\.zliwości jakiegokolwiek dyskursu, nie jego przedmiotem
    (uw.~\ref{rem:0-horyzont}).
\end{itemize}

% -------------------------------------------------------------------
\subsection{Dowody}%
\label{app:0-dowody}
% -------------------------------------------------------------------

\begin{proof}[Dow\'od prop.~\ref{prop:0-endofunktor}]
(i) Endofunktor $\neg: \mathbf{0} \to \mathbf{0}$ jest funkcj\k{a}
pust\k{a} i~zachodzi $\neg = \mathrm{id}_{\mathbf{0}}$; warunek
$\neg \neq \mathrm{id}$ sprzeczny. (ii) Rozwa\.zmy kategori\k{e}
dyskretn\k{a} $\mathcal{D}_n$ z~$n \geq 1$ obiektami. Endofunktor
$\neg$ musi by\'c permutacj\k{a} obiekt\'ow z~$\neg^2 = \mathrm{id}$
i~$\neg \neq \mathrm{id}$, czyli inwolucj\k{a} bez punkt\'ow
sta\l{}ych. Dla $n=1$: brak takiej permutacji. Dla $n=2$: jedyna
mo\.zliwo\'s\'c $+ \leftrightarrow -$. (iii) Automorfizmy
$\mathcal{D}_2$ zachowuj\k{a}ce $\neg$ to $\{\mathrm{id}, \neg\}
\cong \Zp$. \qed
\end{proof}

\begin{proof}[Szkic dow.~prop.~\ref{prop:0-clifford}]
(i) Brak generator\'ow: $\mathrm{Cl}(0,0) = \mathbb{R}$ bez $\Zp$-gradacji.
(ii) Jeden generator $e$, $e^2 = +1$: $\mathrm{Cl}(1,0) = \mathbb{R}[e]/(e^2-1)
\cong \mathbb{R} \oplus \mathbb{R}$ z~projekcjami $(1 \pm e)/2$.
Gradacja $\Zp$: cz\k{e}\'s\'c parzysta $\mathbb{R}\cdot 1$, nieparzysta
$\mathbb{R}\cdot e$. (iii) Alternatywne gradacje $\mathbb{Z}_n$ dla
$n \geq 3$ wymagaj\k{a} pierwiastk\'ow jedno\'sci w~ciele, czyli
uprzedniego rozszerzenia do $\mathbb{C}$ --- sprzeczne z~minimalno\'sci\k{a}.
\qed
\end{proof}

\begin{proof}[Dow\'od prop.~\ref{prop:0-niestabilnosc}]
(i) Trywialne: $\mathrm{id}(x) = x$ dla ka\.zdego $x$. (ii)
Obserwacyjna r\'ownowa\.zno\'s\'c wynika z~braku struktury wyr\'o\.zniaj\k{a}cej
--- dwa stany r\'o\.znice si\k{e} tylko ,,etykietami'' niewidocznymi
spoza przestrzeni stan\'ow. (iii) Dla $\sigma \neq \mathrm{id}$
z~$\sigma^2 = \mathrm{id}$, ka\.zda orbita ma rozmiar 1 (punkty
sta\l{}e $\sigma$) lub 2. Odrzucaj\k{a}c trywialny przypadek wszystkich
punkt\'ow sta\l{}ych ($\sigma = \mathrm{id}$), istnieje co najmniej
jedna orbita $\{x, \sigma x\}$ rozmiaru~2, generuj\k{a}ca $\Zp$. \qed
\end{proof}

\begin{proof}[Szkic dow.~tw.~\ref{thm:0-rownowaznosc}]
Ka\.zde z~trzech sformu\l{}owa\'n produkuje minimaln\k{a} struktur\k{e}
$(\mathcal{T}_{\min}, \sigma)$ z~$|\mathcal{T}_{\min}| = 2$,
$\sigma^2 = \mathrm{id}$, $\sigma$ bez punkt\'ow sta\l{}ych.
Przez jednoznaczno\'s\'c grupy rz\k{e}du 2 wszystkie takie struktury
s\k{a} izomorficzne. Szczeg\'o\l{}y kategoryjne: tw.~Cayleya dla
grup sko\'nczonych + funktor zapominaj\k{a}cy kategoria\,$\to$\,grupa
automorfizm\'ow. \qed
\end{proof}

% -------------------------------------------------------------------
\subsection{Status epistemiczny i~propozycja publikacyjna}%
\label{app:0-status}
% -------------------------------------------------------------------

\begin{itemize}[nosep]
  \item Aksjomat~0 nie jest twierdzeniem empirycznym; jest
    \textbf{warunkiem mo\.zliwo\'sci} sformu\l{}owania TGP.
  \item Nie jest te\.z ,,wolnym parametrem'' --- nie mo\.zna go
    zast\k{a}pi\'c inn\k{a} grup\k{a} bez utraty minimalno\'sci.
  \item Wszystkie predykcje numeryczne TGP pozostaj\k{a} bez zmian;
    aksjomat~0 jedynie \textbf{reorganizuje fundament} logiczny.
  \item Propozycja publikacyjna: w\l{}\k{a}czy\'c niniejszy dodatek
    do pe\l{}nego manuskryptu jako \emph{Dodatek~0} (przed
    dod.~\ref{app:notacja}), a~w~sek.~\ref{sec:ontologia} doda\'c
    jednoakapitow\k{a} wzmiank\k{e} z~odnośnikiem.
\end{itemize}

\begin{remark}[Filozoficzny kontekst]\label{rem:0-filozofia}
Konstrukcja ta jest zbie\.zna z~programami Peirce'a (,,distinction
as primitive''), Wheelera (,,it from bit''), Spencera-Browna
(,,laws of form''), a~z~filozoficznej strony --- z~Hegla
dialektyk\k{a} bycia i nicości. Nie jest to jednak argument
filozoficzny: trzy sformu\l{}owania prop.~\ref{prop:0-endofunktor},
\ref{prop:0-clifford} i~\ref{prop:0-niestabilnosc} s\k{a}
matematycznie \'scis\l{}e i~weryfikowalne w~ramach teorii kategorii,
algebry Clifforda i~dynamiki przestrzeni stan\'ow odpowiednio.
\end{remark}

\begin{observation}[Linia otwarta --- przysz\l{}a praca]%
\label{obs:0-linia}
Trzy r\'ownowa\.zne sformu\l{}owania sugeruj\k{a} istnienie
\textbf{czwartego} --- topologicznego: \'sladu $\pi_0$ na kategorii
przestrzeni topologicznych. Hipoteza: ,,nicość topologiczna'' (zbi\'or
pusty) ma $\pi_0 = 0$, a~minimalny obiekt o~$\pi_0 = \Zp$ to
$S^0 = \{+, -\}$ --- zero-sfera. To mog\l{}oby po\l{}\k{a}czy\'c
aksjomat~0 z~topologi\k{a} solitonu (dod.~\ref{app:E-pi1}),
dodatkowo domykaj\k{a}c struktur\k{e} przez $\pi_1 = \Zp$ ju\.z
obecn\k{a} w~TGP.
\end{observation}

% ===================================================================
%  Koniec dodatku 0
% ===================================================================
